% © 2026 Ing. David Jaroš — CC BY-NC-ND 4.0
%
% This work is licensed under a Creative Commons Attribution-NonCommercial-NoDerivatives
% 4.0 International License (CC BY-NC-ND 4.0).
%
% License History: Earlier drafts (up to v0.3) were released under CC BY 4.0.
% From v0.4 onward, all material is released under CC BY-NC-ND 4.0 to protect
% the integrity of the theoretical work during ongoing academic development.

\documentclass[11pt,a4paper]{article}
\usepackage{amsmath, amssymb, amsthm}
\usepackage{hyperref}
\usepackage{geometry}
\geometry{margin=2.5cm}

\newtheorem{theorem}{Theorem}
\newtheorem{proposition}{Proposition}
\newtheorem{definition}{Definition}
\newtheorem{remark}{Remark}
\newtheorem{conjecture}{Conjecture}

\title{Hecke Eigenvalue Triples and Lepton Mass Ratios\\
\large Arithmetic Structure, Modular Forms, and Spectral Attractors in UBT}
\author{Ing.\ David Jaroš}
\date{2026}

\begin{document}
\maketitle

\begin{abstract}
We analyse the arithmetic structure of the Hecke eigenvalue triples that
numerically reproduce the charged lepton mass ratios in Unified Biquaternion
Theory (UBT).  The observation that Set~A (at prime $p = 137$) and Set~B
(at $p = 139$) of modular newforms reproduce $m_\mu/m_e$ and $m_\tau/m_e$
to 0.02\% and 0.10\% accuracy is documented, and the theoretical mechanisms
that could explain this coincidence are examined.  The focus is on the
arithmetic structure rather than numerical search.
All claims are labelled with their proof status.
\end{abstract}

\tableofcontents

% ======================================================================
\section{Recap of the Numerical Observation}
\label{sec:data}
% ======================================================================

The following is a summary of the results documented in
\texttt{research/mirror\_sector\_modular\_status.md} and verified by
SageMath computation (\texttt{research\_tracks/three\_generations/step5\_hecke\_search\_results.tex}).

\subsection{Set A: Hecke eigenvalues at $p = 137$}

\begin{center}
\begin{tabular}{lccccc}
\hline
Form & Level $N$ & Weight $k$ & $a_{137}$ & Ratio & Target & Error \\
\hline
$f_1$ (electron ref.) & 76 & 2 & $-11$ & 1 & 1 & — \\
$f_2$ (muon) & 7 & 4 & $+2274$ & 206.727 & 206.768 & 0.02\% \\
$f_3$ (tau) & 208 & 6 & $-38286$ & 3480.5 & 3477.23 & 0.10\% \\
\hline
\end{tabular}
\end{center}

\textbf{Status:} Numerical observation \textbf{[NO]}.  Independently
reproducible with SageMath.

\subsection{Lepton mass ratios (experimental)}

\begin{equation}
  \frac{m_\mu}{m_e} = 206.7682831(46), \qquad
  \frac{m_\tau}{m_e} = 3477.23(23).
\end{equation}

\begin{remark}[Epistemic status]
  All mass-ratio agreements cited in this document are \textbf{numerical
  observations at specific primes}.  No theoretical derivation connecting
  the Hecke eigenvalues to physical lepton masses is known.  The results
  are intriguing but must not be presented as predictions of UBT without
  a theoretical mechanism.
\end{remark}

% ======================================================================
\section{Arithmetic Structure of the Modular Forms}
\label{sec:arith}
% ======================================================================

\subsection{Level factorizations}

The levels of the three Set~A forms are:
\begin{align}
  N_1 &= 76 = 4 \times 19, \\
  N_2 &= 7 \quad \text{(prime)}, \\
  N_3 &= 208 = 16 \times 13.
\end{align}

\begin{remark}[Prime structure]
  The level $N_2 = 7$ is the unique minimal prime level for a weight-4 newform
  with the observed $a_{137}$ value.  The occurrence of $N_3 = 208 = 16 \times 13$
  is structurally non-trivial; the factor of 16 suggests a connection to the
  2-adic structure.  Status: \textbf{[NO]} — observed, no theoretical explanation.
\end{remark}

\subsection{Weight progression}

The weights $k = 2, 4, 6$ form a regular progression with step 2.

\begin{remark}[Weight-2 step]
  In the theory of modular forms, even weights $k = 2, 4, 6, \ldots$ are
  precisely those occurring for modular forms on $\Gamma_0(N)$ with trivial
  character.  The step of 2 is consistent with the three generations arising
  from three even winding modes (see \texttt{research/generation\_structure.tex}
  §\ref{sec:modular_arithmetic}).  Status: \textbf{[L2]}.
\end{remark}

\subsection{$L$-function structure}

The $L$-function of each form $f_k$ of weight $k$ is:
\begin{equation}
  L(f_k, s) = \sum_{n=1}^\infty \frac{a_n(f_k)}{n^s}
  = \prod_p \frac{1}{1 - a_p p^{-s} + \chi(p) p^{k-1-2s}},
\end{equation}
where $\chi$ is the nebentypus character (trivial here).

\begin{remark}[GRH and zeros]
  Under the Generalized Riemann Hypothesis, all non-trivial zeros of $L(f_k, s)$
  lie on the critical line $\mathrm{Re}(s) = k/2$.  The critical line varies
  with weight $k$.  If the lepton masses are related to the zeros of the
  $L$-functions, the weight progression $k = 2, 4, 6$ would produce three
  distinct critical lines, consistent with three mass scales.
  Status: \textbf{[L2]} — speculative structural observation.
\end{remark}

% ======================================================================
\section{Stability of the $p = 137$ and $p = 139$ Sectors}
\label{sec:stability}
% ======================================================================

\subsection{The twin-prime pair $(137, 139)$}

The primes $p = 137$ and $p = 139$ form a twin-prime pair separated by 2.

\begin{remark}[Why twin primes?]
  The appearance of two adjacent primes is not explained by any known
  theorem in modular form theory.  It may be a numerical coincidence,
  or it may reflect the parity symmetry $n \leftrightarrow -n$ of the
  KK tower (which maps winding mode $+n$ to $-n$).
  Status: \textbf{[L2]}.
\end{remark}

\subsection{Stability check: nearby primes}

\begin{remark}[Sensitivity to prime choice]
  The Hecke eigenvalue ratios are specific to $p = 137$ and do not hold
  at nearby primes such as $p = 131$ or $p = 149$.  This sensitivity is
  documented in \texttt{reports/hecke\_lepton/statistical\_significance.md}.

  The statistical significance of the match at $p = 137$ is approximately
  3.5$\sigma$ above random chance for a triple of random modular forms.
  This is significant but not conclusive.  Status: \textbf{[NO]}.
\end{remark}

% ======================================================================
\section{Spectral Attractors}
\label{sec:attractors}
% ======================================================================

\subsection{Definition}

\begin{definition}[Spectral attractor]\label{def:attractor}
  A \emph{spectral attractor} at prime $p$ is a modular form $f$ of level $N$
  and weight $k$ such that the Hecke eigenvalue $a_p(f)$ is an extremum
  (maximum or minimum) over all forms in the family $\{f_{N, k}\}_{N \leq N_{\max}}$.
\end{definition}

The hypothesis is that the Set~A forms at $p = 137$ are spectral attractors in
the sense that their Hecke eigenvalues are determined by an extremisation
principle related to the UBT effective potential.

\begin{conjecture}[Hecke eigenvalues from spectral attractor]\label{conj:attractor}
  The three modular forms in Set~A are spectral attractors of the Hecke operator
  $T_{137}$ acting on forms of weights $k = 2, 4, 6$, in the sense that they
  have Hecke eigenvalues closest to the golden ratio times the Ramanujan bound
  $|a_{137}| \leq 2 \cdot 137^{(k-1)/2}$.

  \textbf{Status: [L2]} — not verified; the Ramanujan bound is saturated at
  Hecke eigenvalues $|a_{137}|^2 = 4 \cdot 137^{k-1}$, not at the lepton
  mass ratios.  This remains a speculative direction.
\end{conjecture}

% ======================================================================
\section{Relation to UBT Field Dynamics}
\label{sec:dynamics}
% ======================================================================

\subsection{Modular forms from the KK partition function}

The KK partition function $Z_{\mathrm{KK}}(\tau) = \vartheta_3^3(\tau)$ is
a modular form of weight $3/2$ under a congruence subgroup.  Its Hecke
eigenvalues at prime $p$ are:
\begin{equation}
  T_p\, Z_{\mathrm{KK}} = (1 + p^{3/2})\, Z_{\mathrm{KK}} + \ldots
\end{equation}

\begin{remark}[Hecke weight $3/2$ vs integer weight]
  The KK partition function is a half-integer weight form (weight $3/2$),
  while the Set~A forms have integer weights $k = 2, 4, 6$.  These are
  mathematically distinct spaces.  A connection between them would require
  a Shimura correspondence or Waldspurger-type lift.
  Status: \textbf{[O]} — a theoretical mechanism is completely unknown.
\end{remark}

\subsection{Physical interpretation candidates}

Two candidate physical mechanisms connecting Hecke eigenvalues to masses are:

\begin{enumerate}
  \item \textbf{Period map}: Fermion masses are periods of differential forms
    on arithmetic varieties.  The Hecke eigenvalues label the arithmetic
    variety (modular curve) whose period reproduces the mass.
    Status: \textbf{[L2]} — inspired by the theory of complex multiplication.

  \item \textbf{Spectral triple}: The Dirac operator on the noncommutative
    geometry of the SM spectral triple has a spectrum related to Hecke
    eigenvalues.  Status: \textbf{[L2]} — inspired by Connes NCG.
\end{enumerate}

% ======================================================================
\section{Open Problems}
\label{sec:open}
% ======================================================================

\begin{enumerate}
  \item \textbf{[O] Theoretical derivation}: No mechanism relating Hecke
    eigenvalues at $p = 137$ to lepton masses is known.  This is the
    central open problem.

  \item \textbf{[O] Why $p = 137$?}: The fine-structure constant
    $\alpha \approx 1/137$ is defined in terms of $p = 137$.  The
    appearance of the same prime in both the $\alpha$ formula and the
    Hecke observation may be a coincidence or a deep structural feature.

  \item \textbf{[O] Set~B consistency}: The twin-prime mirror sector at
    $p = 139$ (Set~B) reproduces the same ratios.  The theoretical reason
    for this mirror symmetry is unknown.  See
    \texttt{research/mirror\_sector\_dynamics.tex}.
\end{enumerate}

\end{document}
